% Generated draft: requires mathematical and source review.
\documentclass[11pt]{article}
\usepackage[margin=25mm]{geometry}
\usepackage{fontspec}
\setmainfont{Latin Modern Roman}
\usepackage{amsmath,amssymb,mathtools,unicode-math}
\setmathfont{Latin Modern Math}
\usepackage{enumitem,needspace,xurl,hyperref}
\setlength{\parindent}{0pt}
\setlength{\parskip}{5pt}
\setlength{\emergencystretch}{3em}
\newcommand{\N}{\mathbb{N}}
\newcommand{\Z}{\mathbb{Z}}
\newcommand{\Q}{\mathbb{Q}}
\newcommand{\R}{\mathbb{R}}
\newcommand{\C}{\mathbb{C}}
\newcommand{\F}{\mathbb{F}}
\newcommand{\A}{\mathbb{A}}
\newcommand{\PP}{\mathbb{P}}
\newcommand{\E}{\mathbb{E}}
\newcommand{\eps}{\varepsilon}
\newcommand{\dd}{\,\mathrm{d}}
\begin{document}
\section*{UnsolvedMath ID: 20001944}
\textbf{Source problem code:} catalogue entry
\subsection{Definitions and mathematical statement}
The supplied canonical Question 3, attributed to Cohn--Miller, asks the following three questions about their determinant construction:
\begin{enumerate}
\item[(A)] Can the determinants be rescaled as prescribed in their Conjecture 3.2?
\item[(B)] Do the resulting rescaled determinants converge uniformly on compact sets to a limiting function?
\item[(C)] Does that limit have any sign changes other than the forced ones?
\end{enumerate}
The sought combination of answers is yes to (A), yes to (B), and no to (C). The final sentence, beginning with the consequence of that combination, is truncated after ``then the answer to''. The determinant matrices, their arguments and domains, Conjecture 3.2, the rescaling, the forced sign changes, and the target of this final implication are absent. Consequently these questions cannot be made into a fully specified determinant-sequence assertion from the supplied canonical excerpt.

The research aid distinguishes the literal 2004 reference from a precise public 2016 polynomial--Gaussian moving-root construction. It reports that the pre-2004 draft and its Conjecture 3.2 were not identified, and makes no claim about the exact historical normalization. For the public construction it reports finite existence and uniqueness, while convergence of the specified approximants and exclusion of extra real zeros of their hypothetical limit were not found proved. These statements concern that public sequence, not a verified identification with the historical one.

The aid's finite determinant criterion uses quantities \(A_0,A_1,B_0,B_1\) and states that the normalized linear combination satisfying the remaining Fourier-side double-zero condition exists uniquely exactly when
\[
 \Delta_k=B_0A_1+B_1A_0\ne0.
\]
It gives the two coefficients as \(B_1/\Delta_k\) and \(B_0/\Delta_k\). The underlying Hermite determinants, the definitions and \(k\)-dependence of these four quantities, the basis/sign conventions for the combination, and the location of the double zero are not supplied. The displayed criterion is thus a recorded formula with unresolved inputs, not a stand-alone theorem about arbitrary determinants.

Its conditional limiting claim concerns the specified public approximants in dimension \(d=8\) or \(d=24\). To make the analytic hypotheses explicit without inventing the missing construction, write \(f_k\) for those radial functions on \(\R^d\). The reported hypotheses include local uniform convergence \(f_k\to f\), a common envelope
\[
 |f_k(x)|\le H(x)\quad\text{for all }k\text{ and almost every }x,
 \qquad H\in L^1(\R^d),
\]
and a radial Schwartz limit \(f\). Together with the construction's interpolation and normalization data, the aid claims these force \(f\) to be the known dimension-\(d\) magic function by interpolation uniqueness. Those data and the normalization of ``magic function'' are not reproduced, so this cannot be read as a uniqueness theorem for arbitrary radial Schwartz limits. The aid further reports full-sequence convergence from relative compactness and the stated tail/limit hypotheses; this requires every subsequential limit to fall within the same uniqueness class, not merely one chosen limit.

The aid says the original optimal-function question was solved independently using modular forms and interpolation methods. That reported resolution does not itself prove convergence of this particular determinant sequence. It also discusses absence of extra real zeros, whereas the historical question (C) asks about extra sign changes; a zero need not change sign. The precise equivalence intended here is not supplied.
\subsection{Short English statement}
The historical problem asks for a determinant normalization, compact-uniform convergence, and control of the limit's sign changes. The excerpt does not define the determinants and ends mid-sentence. The aid reports a finite criterion and a conditional limit-identification result for a later public construction, while retaining the convergence problem and the historical identification as unresolved.
\subsection{Sources}
\begin{itemize}
\item Local UnsolvedMath registry, permanent ID 20001944, exact source code catalogue entry; canonical \texttt{aim-geometry-notes.json}, \texttt{notes[281]}, tag \texttt{question}; AIM Workshop Problem Lists, Geometry, workshop \emph{Sphere packings, lattices, and infinite dimensional algebra}, Question 3. Created and updated 14 August 2026.
\item Supplied AIM source URL: \url{https://aimath.org/WWN/spherepacking/spherepacking.pdf}. The supplied verification explicitly identifies the canonical record as truncated. The PDF was not accessed.
\item Cohn--Miller Conjecture 3.2 is cited without its text or a bibliographic entry. The aid distinguishes an unidentified pre-2004 draft from a public 2016 construction, but supplies no title, publication entry, determinant definitions, or authoritative theorem references for that construction or for the modular-form/interpolation resolution. These source gaps require review.
\item Supplied attempt 1: \texttt{partial\_result}, problem status at run \texttt{partially\_solved}, rigor label \texttt{proved}, novelty confidence \texttt{low}; classification \texttt{PARTIAL-PROGRESS}, registry status \texttt{partially\_solved}. Candidate: determinant nondegeneracy criterion and controlled-limit rigidity theorem. Statements that convergence and zero exclusion were not found proved are the aid's dated search assessment, not a new literature conclusion.
\item Ulam AI / UnsolvedMath Contributors; CC BY 4.0. Machine-generated research aid, not peer reviewed; independently verify before citation.
\end{itemize}
\end{document}
